\chapter{Architectural and Theoretical Background}

This chapter presents the main architectural principles and technologies
used in the design of Football OS. The objective is to explain the
general technical approach behind the system and provide a clear
understanding of how its components interact at a high level.

% -------------------------------------------------------
\section{Architectural Foundations}

Football OS is designed as a distributed backend platform composed of
independent services supported by a shared SDK. The system is organized
around four main components:

\begin{itemize}
    \item \textbf{API Gateway}: handles authentication, request enrichment, and routing,
    \item \textbf{Governance Service}: manages identity, authorization policies, and audit logs,
    \item \textbf{Workspace Service}: manages operational features such as documents, events, and notifications,
    \item \textbf{Shared SDK}: provides reusable modules for cross-cutting concerns.
\end{itemize}

All client requests pass through the API Gateway before reaching backend
services, ensuring consistent security, validation, and routing across
the platform.
\begin{figure}[H]
    \centering
    \includegraphics[width=\textwidth]{images/high-level-architecture.png}
    \caption{Football Os - high level architecture }
    \label{fig:high-level-architecture}
\end{figure}

\subsection{Service-Oriented Architecture}

Football OS follows a service-oriented architecture where each service
owns its data and business logic. This separation allows components to
evolve independently and reduces coupling between different domains.

The governance service is responsible for core system concerns such as
identity and authorization, while the workspace service focuses on
user-facing operations such as event management and document handling.
This separation ensures that changes in one domain do not directly impact
the others.

\subsection{Event-Driven Architecture}

To improve scalability and decoupling, Football OS adopts an event-driven
architecture based on Apache Kafka.

Instead of relying only on direct service calls, services communicate
through events. For example, when an event or document is created in the
workspace service, a corresponding event is published and can be consumed
by other parts of the system, such as audit or notification components.

This approach allows asynchronous processing, reduces dependencies
between services, and makes the system more resilient to failures.

\subsection{Layered Architecture}

The system also follows a layered architectural structure to separate
business logic from infrastructure concerns.

At the core, domain logic defines the system's rules and entities.
Above it, application services orchestrate workflows. External concerns
such as databases, messaging systems, and APIs are handled in outer layers.

This structure improves maintainability, testability, and clarity by
ensuring that core logic remains independent of technical details.
\begin{figure}[H]
    \centering
    \includegraphics[width=\textwidth]{images/onion architecture.png}
    \caption{Football Os - Layered Architecture }
    \label{fig:high-level-architecture}
\end{figure}
% -------------------------------------------------------
\section{Shared SDK and Design Principles}

A key architectural decision in Football OS is the use of a shared backend
SDK (\texttt{fos-sdk}) across all services.

The SDK provides reusable modules that handle common concerns such as:
\begin{itemize}
    \item Event communication,
    \item Security context handling,
    \item Policy evaluation integration,
    \item Canonical data resolution,
    \item Storage abstraction.
\end{itemize}

By centralizing these functionalities, the SDK avoids duplication and
ensures consistent behavior across services. This allows each service to
focus only on its core business logic while relying on standardized
components for cross-cutting concerns.

The SDK is built using well-known design principles such as abstraction,
modularity, and separation of concerns, which simplify development and
improve system coherence.

% -------------------------------------------------------
\section{Security Model}

\subsection{Authentication}

Authentication in Football OS is handled using Keycloak, which provides
user management and issues JSON Web Tokens (JWT). Each request includes
a token that identifies the user and their roles.

The API Gateway validates this token before forwarding the request,
ensuring that only authenticated users can access the system.

\subsection{Authorization}

Football OS uses an Attribute-Based Access Control (ABAC) model to
control access to resources.

Access decisions are based on multiple factors, including:
\begin{itemize}
    \item The user's role,
    \item The requested action,
    \item The type and state of the resource.
\end{itemize}

This allows fine-grained control over permissions, especially for
sensitive data such as medical documents.

\subsection{Policy Evaluation}

Authorization logic is externalized using Open Policy Agent (OPA).
Services delegate access decisions to the governance service, which
evaluates policies and returns an allow or deny response.

This approach centralizes security rules and allows them to be updated
without modifying application code.

% -------------------------------------------------------
\section{Communication Model}

Football OS uses two complementary communication mechanisms:

\begin{itemize}
    \item \textbf{REST APIs}: for synchronous communication between clients and services,
    \item \textbf{Kafka}: for asynchronous communication between services.
\end{itemize}

All external requests go through the API Gateway, which ensures proper
authentication and routing. Internally, services may use direct calls
for immediate decisions and events for background processing.

This hybrid approach balances responsiveness and scalability.

% -------------------------------------------------------
\section{Data Management}

Football OS uses a hybrid data management strategy to match different
types of data.

\subsection{MongoDB}

MongoDB is used in the workspace service to store documents, events,
and notifications. These data structures are often complex and nested,
making a document-oriented database a suitable choice.

\subsection{PostgreSQL}

PostgreSQL is used in the governance service to manage structured data
such as users, teams, and audit logs. It provides strong consistency
and supports relational data modeling.

\subsection{Redis}

Redis is used for caching and rate limiting. It improves performance by
reducing repeated computations and controlling request traffic at the
gateway level.

\subsection{Canonical Reference Approach}

To ensure data consistency across services, Football OS avoids duplicating
data. Instead, services store references to external entities and resolve
them when needed through the governance service.

% -------------------------------------------------------
\section{Infrastructure}

\subsection{API Gateway}

The API Gateway acts as the single entry point to the system. It handles
authentication, request validation, rate limiting, and routing to backend
services.

\subsection{Containerized Environment}

Docker Compose is used to run the system's infrastructure components,
including databases, messaging systems, and authentication services.
This ensures a consistent and reproducible development environment.

\subsection{Object Storage}

MinIO is used as an object storage system for managing uploaded files.
It provides an S3-compatible interface and is accessed through an
abstraction layer, allowing flexibility in choosing storage providers.

% -------------------------------------------------------
\section*{Conclusion}
\addcontentsline{toc}{section}{Conclusion}

This chapter introduced the main architectural concepts and technologies
used in Football OS. It presented the system as a modular, service-based
platform supported by a shared SDK and combining both synchronous and
asynchronous communication.

These concepts provide the foundation for understanding the system
requirements and detailed design presented in the following chapters.
